\subsection{Other Implementation Details}
\label{subsec:Implementation}

We implemented \toolname\ as a prototype DBMS testing system in Python, consisting of \textcolor{red}{13,455} lines of code.
To support mutation at the structural level, \toolname\ leverages \textsc{SQLGlot}~\cite{tobymao_sqlglot} to parse original DML statements into ASTs, enabling precise application of approximate mutators while preserving syntactic and semantic correctness.
During testing, \toolname\ executes generated statements on target DBMSs and collects execution results for further analysis.
To improve the usability of detected bugs, we adopt \textsc{SQLess}~\cite{lin2024sqless} to simplify failing cases and isolate their root causes, and further apply a deduplication strategy following \textsc{Pinolo}~\cite{hao2023pinolo} to eliminate redundant bug reports.
\toolname\ is primarily designed for MySQL-compatible systems, including TiDB and OceanBase, but its modular architecture allows for straightforward adaptation to other DBMS dialects by customizing syntax- and type-related components.
In practice, such adaptation requires only lightweight modifications, and can be further automated using tools such as \textsc{QTRAN}~\cite{lin2025qtran}, which facilitates cross-dialect deployment of metamorphic testing techniques.